\documentclass[12pt, letterpaper]{article}

%-------Packages---------
\usepackage{newpxtext,newpxmath} % Palatino-like text & math (modern)
\usepackage{bboldx}
\usepackage{mathtools}
\usepackage{tikz}
\usetikzlibrary{calc,intersections}
\usepackage{tikz-cd}
\usepackage[all,arc]{xy}
\usepackage{graphicx}

\usepackage{enumerate}
\usepackage[dvipsnames]{xcolor}
\usepackage{float}
\usepackage[left=1in,top=1in,right=1in,bottom=1in]{geometry}
\usepackage{fancyhdr}
\usepackage{multicol}
\usepackage{setspace}

\usepackage{dsfont}
\usepackage{mathrsfs}

\usepackage{tcolorbox}
\tcbuselibrary{theorems,skins,breakable}

\usepackage{xparse}
\usepackage{import}
\usepackage[utf8]{inputenc}

\usepackage{hyperref}
\usepackage{cleveref}

\usepackage{xcolor, ifthen, tcolorbox}
\tcbuselibrary{theorems,skins,breakable}
% Theorem System
% The following environments are provided:
%   New Problem:        \begin{problem} ...
%   New Question Header:   \qh (then followed by subproblems)
%   Subproblem:     \begin{subproblem} ...
%   Problem Proof:  \begin{problemproof} ...
%   Proof:          \\begin{pf}{color} ...
%   Remark:         \rmk (sentence), \rmkb (block)

% Define custom colors
\definecolor{thmscol}{HTML}{F4565B} %For Problems
\definecolor{lemscol}{HTML}{FE844C} %For Lemmes
\definecolor{prpscol}{HTML}{00A7EF} %For proposition
\definecolor{corscol}{HTML}{23DAFF} %For corrolaries
\definecolor{rmkscol}{HTML}{6DEEC5} %For remarks
\definecolor{defscol}{HTML}{18C991} %For definitions
\definecolor{exmscol}{HTML}{7DB800} %For examples
\definecolor{clmscol}{HTML}{165c58} %For claims

%Change between dark(black)/light(white) mode
\newcommand{\colormode}{white}
\ifthenelse{\equal{\colormode}{white}}
      {\newcommand{\TextColor}{black}}
      {\newcommand{\TextColor}{white}}
\newcommand{\pbcolormain}{thmscol!80!\colormode}
\newcommand{\pbcolorsecond}{thmscol!38!\colormode}


% ============================
% Problem Counter
% ============================
\newcounter{subproblemcounter}
\newcounter{problemcounter}

\newcommand{\q}{
    \setcounter{subproblemcounter}{0}%
    \stepcounter{problemcounter} % Increment problem counter
    \color{\TextColor}Problem \arabic{problemcounter}: % Display the problem counter
}

\newcommand{\stepsub}{
    \stepcounter{subproblemcounter}%
    \color{\TextColor}(\alph{subproblemcounter})
}
% ============================
% Problem
% ============================
\newtcolorbox{pb}[1]{
  enhanced,
  frame hidden,
  titlerule=0mm,
  toptitle=1mm,
  bottomtitle=1mm,
  fonttitle=\bfseries\large,
  coltitle=black,
  colbacktitle=\pbcolormain,
  colback=\pbcolorsecond,
  title={\q #1},
  coltext=\TextColor
}

\newtcolorbox{spb}[1]{
  enhanced,
  frame hidden,
  titlerule=0mm,
  toptitle=1mm,
  bottomtitle=1mm,
  fonttitle=\bfseries\large,
  coltitle=black,
  colbacktitle=\pbcolormain,
  colback=\pbcolorsecond,
  title={\stepsub #1},
  coltext=\TextColor,
  attach boxed title to top left={xshift=3mm,yshift=-2mm},
  boxed title style={
    colback=\pbcolormain,
    colframe=\pbcolormain,
  }
}

\newtcolorbox{pbh}[1]{
    enhanced,
    frame hidden,
    titlerule=0mm,
    toptitle=1mm,
    bottomtitle=1mm,
    fonttitle=\bfseries\large,
    coltitle=black,
    colbacktitle=\pbcolormain,
    colback=\pbcolorsecond,
    title={\q #1},
    boxrule=0pt,
    interior hidden,
    attach boxed title to top left={xshift=3mm,yshift=-2mm},
    boxed title style={
        colback=\pbcolormain,
        colframe=\pbcolormain,
    }
}

\newenvironment{problem}[1]{%
  \newpage
  \begin{pb}{#1}
    }{
  \end{pb}
}

\newenvironment{subproblem}[1]{%
  \begin{spb}{#1}
    }{
  \end{spb}
}

\newenvironment{problemproof}{%
    \begin{pf}{\pbcolormain}
        }{
    \end{pf}
}

\newcommand{\qh}[1][]{
    \newpage
    \begin{pbh}{#1}\end{pbh} % Display the problem counter
}

% ============================
% Claim
% ============================

\newtcbtheorem[number within=problemcounter]{claim}{Claim}
{
    enhanced,
    frame hidden,
    titlerule=0mm,
    toptitle=1mm,
    bottomtitle=1mm,
    fonttitle=\bfseries\large,
    coltitle=black,
    colbacktitle=clmscol!50!\colormode,
    colback=clmscol!20!\colormode,
}{clm}

\NewDocumentCommand{\clm}{m+m}{
    \begin{claim*}{#1}{}
        #2
    \end{claim*}
}

\newenvironment{clmpf}{
	{\noindent{\it \textbf{Proof.}}
	\setlength{\parindent}{0pt}
	}
	\tcolorbox[blanker,breakable,left=5mm,parbox=false,
    before upper={\parindent15pt},
    after skip=10pt,
	borderline west={1mm}{0pt}{clmscol!50!\colormode}]
}{
    \textcolor{clmscol!50!\colormode}{\hbox{}\nobreak\hfill$\blacksquare$}
    \endtcolorbox
}

\NewDocumentCommand{\clmp}{m+m+m}{
    \begin{claim*}{#1}{}
        #2
    \end{claim*}

    \begin{clmpf}
        #3
    \end{clmpf}
}


% ============================
% Proof
% ============================

\newenvironment{pf}[1]{%
  \def\pfcolor{#1}% Store the color in a macro
  \noindent{\it \textbf{Proof.}}%
  \tcolorbox[blanker,breakable,left=5mm,parbox=false,%
    before upper={\parindent15pt},%
    after skip=10pt,%
    borderline west={1mm}{0pt}{\pfcolor},%
    coltext=\TextColor
  ]%
  \setlength{\parindent}{0pt}%
}{%
  \textcolor{\pfcolor}{\hbox{}\nobreak\hfill$\blacksquare$}%
  \endtcolorbox%
}

\newtcolorbox{solution}{
  blanker,
  breakable,
  left=5mm,
  parbox=false,%
  before upper={\parindent15pt},%
  after skip=10pt,%
  borderline west={1mm}{0pt}{\pbcolormain},%
  coltext=\TextColor
}


% ============================
% Remark
% ============================


\NewDocumentCommand{\rmk}{+m}{
    {\it \color{rmkscol!100!white}#1}
}

\newenvironment{remark}{
    \par
    \vspace{5pt}
    \begin{minipage}{\textwidth}
        {\par\noindent{\textbf{Remark.}}}
        \tcolorbox[blanker,breakable,left=5mm,
        before skip=10pt,after skip=10pt,
        borderline west={1mm}{0pt}{rmkscol!80!\colormode},
        coltext=\TextColor
        ]
}{
        \endtcolorbox
    \end{minipage}
    \vspace{5pt}
}

\NewDocumentCommand{\rmkb}{+m}{
    \begin{remark}
        #1
    \end{remark}
}

\definecolor{rmkscol}{HTML}{F4565B}
\renewcommand{\stepsub}{
    \stepcounter{subproblemcounter}%
    \color{\TextColor}(\arabic{subproblemcounter})
}

% Define a new command for problems
\renewcommand{\headrule}{\color{\TextColor}\hrulefill}
\newcommand{\C}{\mathbb{C}}
\newcommand{\R}{\mathbb{R}}
\newcommand{\Q}{\mathbb{Q}}
\newcommand{\Z}{\mathbb{Z}}
\newcommand{\N}{\mathbb{N}}
\newcommand{\p}{\mathfrak{p}}
\newcommand{\F}{\mathbb{F}}
\newcommand{\contr}{\Rightarrow \Leftarrow}
\newcommand{\s}{\(\,\)}
\renewcommand{\t}[1]{\text{#1}}
\newcommand{\Spec}{\mathrm{Spec}}
\newcommand{\Ass}{\mathrm{Ass}}
\newcommand{\Supp}{\mathrm{Supp}}
\newcommand{\Ann}{\mathrm{Ann}}
\newcommand{\mf}[1]{\mathfrak{#1}}
\newcommand{\codim}{\mathrm{codim}}

% Meta data
\setstretch{1.2}

\begin{document}
\color{\TextColor}
\pagecolor{\colormode}
\setlength{\parindent}{0pt}
\pagestyle{fancy}
\fancyhf{}
\fancyhead[LOE]{\color{\TextColor}\slshape{\textbf{Vincent Ferrara}}}
\fancyhead[ROE]{\color{\TextColor}\jobname --- \thepage}
\fancyhead[C]{\color{\TextColor}Math 493}

\stepcounter{problemcounter}
\begin{problem}{}
    Let $G$ be a group and let $a$ be an element of $G$. Show that $g \mapsto aga^{-1}$ is a group homorphism $G \to G$.
\end{problem}
\begin{problemproof}
    Let $\varphi : G \to G$ s.t. $\varphi(g)=aga^{-1}$

    Let $g,h \in G$
    \begin{align*}
        \varphi(gh)=agha^{-1}=aga^{-1}aha^{-1}=\varphi(g)\varphi(h)
    \end{align*}
\end{problemproof}

\begin{problem}{}
    Let $G$ and $H$ be two groups and let $\alpha: G \to H$ and $\beta: G \to H$ be group homorphisms. For $g \in G$, define $(\alpha\beta)(g)=\alpha(g)*\beta(g)$.
\end{problem}

\begin{subproblem}{}
    If $H$ is commutative, show that $\alpha\beta$ is a group homorphism.
\end{subproblem}
\begin{problemproof}
    Assume $H$ is commutative.

    Let $g,h \in G$.
    \begin{align*}
        (\alpha\beta)(gh) &= \alpha(gh) * \beta(gh)\\ 
                          &= \alpha(g) * \alpha(h) * \beta(g) * \beta(h)\\
                          &= \alpha(g) * \beta(g) * \alpha(h) * \beta(h) = (\alpha\beta)(g) * (\alpha\beta)(h)
    \end{align*}
\end{problemproof}
\begin{subproblem}{}
    Give an example of two groups and two homorphisms such that $\alpha\beta$ is not a homorphism.
\end{subproblem}
\begin{problemproof}
    Let $G,H = S_3$ and $\alpha,\beta = \mathsf{id}$.

    Let $g=(12)$ and $h=(23)$. 
    
    $gh=(123)$
    \[(\alpha\beta)(gh)=(123)(123)=(132) \neq e = (12)(12)(23)(23) = (\alpha\beta)(g)(\alpha\beta)(h)\]
\end{problemproof}

\begin{problem}{}
    Let $G$ and $H$ be two groups. We define the \textbf{product group} $G \times H$ to have, as underlying set, the Cartesian product $G \times H$, with operation
    \[(g_1,h_1)(g_2,h_2) = (g_1g_2,h_1h_2).\]
\end{problem}
\stepcounter{subproblemcounter}
\rmkb{
    Define $\Delta : G \to G \times G$ by $\Delta(g)=(g,g)$, and define $\nabla : G \to G \times G$ by $\nabla(g,h)=gh$. 
    }
\begin{subproblem}{}
    Is $\Delta$ always a group homorphism?
\end{subproblem}
\begin{problemproof}
    Yes, let $g,h \in G$.
    \begin{align*}
        \Delta(gh) &= (gh,gh)\\
                   &= (g,g)(h,h) = \Delta(g)\Delta(h)
    \end{align*}
\end{problemproof}
\begin{subproblem}{}
    Is $\nabla$ always a group homorphism?
\end{subproblem}
\begin{problemproof}
    No, let $G=S_3$ and $g=(12)$, $h=(23)$.
    \begin{align*}
        \nabla((g,g)(h,h)) &= \nabla((gh,gh)) = ghgh = (132) \neq
        e=gghh=\nabla(g,g)\nabla(h,h)
    \end{align*}
\end{problemproof}

\begin{problem}{}
    Let $S_n$ act on $\Q[x_1,x_2,\dots,x_n]$ in the obvious way. Let $\Delta = \prod_{1 \leq i < j \leq n}(x_i - x_j) \in \Q[x_1,x_2,\dots,x_n]$. For $\sigma \in S_n$, let $M(\sigma)$ be the corresponding permutation matrix in $\mathsf{GL}_n(\Q)$.
\end{problem}
\begin{subproblem}{}
    Show that
    \[\frac{\sigma(\Delta)}{\Delta}=(-1)^{\left\lvert \{(i,j) \mid 1 \leq i < j \leq n, \sigma(i) > \sigma(j)\} \right\rvert}=\det M(\sigma)\]
\end{subproblem}
\begin{problemproof}
    \[\sigma(\Delta)=\prod_{1 \leq i < j \leq n}(x_{\sigma(i)}-x_{\sigma(j)})\]
    \[
    x_{\sigma(i)}-x_{\sigma(j)} = \begin{cases}
        x_i-x_j, & \sigma(i) < \sigma(j),\\
        -(x_i-x_j), & \sigma(i)>\sigma(j).
    \end{cases}
    \]
    Therefore,
    \[\sigma(\Delta)=(-1)^{\left\lvert \{(i,j) \mid 1 \leq i < j \leq n, \sigma(i) > \sigma(j)\} \right\rvert}\Delta \implies \frac{\sigma(\Delta)}{\Delta}=(-1)^{\left\lvert \{(i,j) \mid 1 \leq i < j \leq n, \sigma(i) > \sigma(j)\} \right\rvert}\]

    A permutation matrix is made by swapping rows. The determinant is 1 if there is an even number of swaps and $-1$ if there is a odd number.

    Thus, \[\det(M(\sigma))=(-1)^{\left\lvert \{(i,j) \mid 1 \leq i < j \leq n, \sigma(i) > \sigma(j)\} \right\rvert}\]
\end{problemproof}
\begin{subproblem}{}
    We defined $\epsilon(\sigma)$ to be given by any of the formulas. Show that $\epsilon$ is a group homomorphism $S_n \to \{\pm 1\}$.
\end{subproblem}
\begin{problemproof}
    Let $\sigma,\tau \in S_n$.
    \[\epsilon(\sigma\tau)=\det (M(\sigma\tau))=\det(M(\sigma)M(\tau))=\det(M(\sigma))\det(M(\tau))=\epsilon(\sigma)\epsilon(\tau)\]
\end{problemproof}
\rmkb{
    The \textbf{alternating group} $A_n$ is defined as the kernel of $\epsilon$.
}
\begin{subproblem}{}
    Show that $S_n$ is generated by the elements $(1j)$ for $2 \leq j \leq n$.
\end{subproblem}
\begin{problemproof}
    Consider $k$-cycle $\sigma=(a_1a_2\dots a_k)$.

    Consider $(1a_1)(1a_k)(1a_{k-1})\dots(1a_2)(1a_1)$

    Let $i \in \{1,\dots,n\}$.

    If $i = a_j$ for $j \in \{1,\dots,k\}$ then

    \[(1a_1)(1a_k)(1a_{k-1})\dots(1a_2)(1a_1)(a_j)=a_{j+1 \mod k}=\sigma(a_j)\]

    If $i \neq a_j$ for $j \in \{1,\dots,k\}$ then

    \[(1a_1)(1a_k)(1a_{k-1})\dots(1a_2)(1a_1)(i)=i=\sigma(i)\]

    Thus, we can make any $k$ cycle out of $(1j)$ for $2 \leq j \leq n$. Thus, we can make any element in $S_n$.
\end{problemproof}
\begin{subproblem}{}
    Show that $A_n$ is generated by the elements $(12k)$ for $3 \leq k \leq n$.
\end{subproblem}
\begin{problemproof}
    Let $\sigma \in A_n$. From (3) we can write $\sigma=(1a_1)(1a_2)\dots(1a_k)$. Since we define $A_n = \ker(\epsilon)$.

    \[\epsilon(\sigma)=\epsilon(1a_1)\epsilon(1a_2)\dots\epsilon(1a_k)=0\]
    Since $\epsilon(1a_i)=-1$ this implies $k$ is even.
    
    This implies that $A_n$ when decomposed has an even amount of transpositions. Thus, we only need to look at $(1a)(1b)$ and generate it with $(12k)$.

    
    If $a = b$:
    $(1a)(1b)=e$

    If $a \neq b$:

    If $a,b \geq 2$:
    $(1a)(1b)=(12a)(12b)(12b)$

    If $a = 2$:
    $(12)(1b)=(12b)(12b)$

    If $b = 2$:
    $(1a)(12)=(12a)$

    In every case, $(1a)(1b)$ is a product of $(12k)$ with $k \geq 3$, and since every $\sigma \in A_n$ can be written with an even amount of transpositions of the form $(1a)$. This implies that $A_n$ is generated by the elements $(12k)$ for $k \geq 3$.
\end{problemproof}

\qh
\begin{subproblem}{}
    Let $n$ be a positive integer and $a$ another integer. Show that there exists $x$ s.t. $ax \equiv 1 \mod n$ if and only if $\gcd(a,n) = 1$.
\end{subproblem}
\begin{problemproof}
    $(\implies)$:

    Assume $\exists x$ s.t. $ax \equiv 1 \mod n$.

    This means there exists a $y \in \Z$ s.t.
    \[ax - 1 = ny \implies ax-ny = 1\]
    Let $d=\gcd(a,n)$. Since $d \mid a$ and $d \mid n$ this implies $d\mid 1$. Thus, $\gcd(a,n) = d = 1$.

    $(\impliedby)$:

    Assume $\gcd(a,n) = 1$.

    By Bezout's Lemma there exists $x,y \in \Z$ s.t. $ax+ny = 1$. Thus $ax - 1 = ny \implies ax \equiv 1 \mod n$.
    Thus, there exists an $x$.
\end{problemproof}

\begin{subproblem}{}
    Let $(\Z \diagup n\Z)^\times$ be the set of residue classes $\bar{a}$ modulo $n$ for which $\gcd(a,n)=1$. Show that $(\Z \diagup n \Z)^{\times}$ is a group.
\end{subproblem}
\begin{problemproof}
    Associativity:

    It is by definition associative.

    Identity element:

    $1 \in (\Z \diagup n\Z)^\times$ and for all $g \in (\Z \diagup n\Z)^\times$ $1g=g=g1$.

    Inverse element:

    Let $g \in (\Z \diagup n\Z)^\times$. This means that $\gcd(g,n)=1$. Thus, by (1) there exists an $x \in \Z$ s.t. $gx \equiv 1 \mod n$. Therefore, $g\bar{x}=1=\bar{x}g$.
    
    Thus, $(\Z \diagup n\Z)^\times$ is a group.
\end{problemproof}

\begin{subproblem}{}
    Let $n$ have prime factorization $n=p_1^{a_1}p_2^{a_2}\dots p_r^{a_r}$. Define a map $(\Z \diagup n\Z)^\times \to \prod (\Z \diagup p_i^{a_i} \Z)^\times$ by $(m \mod \p_1^{a_1}, m \mod p_2^{a_2},\dots,m \mod p_r^{a_r})$, where the product of groups is defined in. Show that this map is an isomorphism of groups.
\end{subproblem}
\begin{problemproof}
    Let $\varphi$ be that map.

    Well-Defined:

    Let $m \equiv m' \mod n$, then since $p_i^{a_i} \mid n$, we have
    \[m \equiv m' \mod p_i^{a_i}\]
    for every i. Hence, each coordinate is well-defined.

    Homomorphism:

    Let $g,h \in (\Z \diagup n\Z)^\times$.
    \begin{align*}
        \varphi(gh) &=(gh \mod p_1^{a_1}, gh \mod p_2^{a_2},\dots,gh\mod p_r^{a_r})\\
                    &=(g \mod p_1^{a_1}, g \mod p_2^{a_2},\dots,g\mod p_r^{a_r})\\ 
                    &\cdot (h \mod p_1^{a_1}, h \mod p_2^{a_2},\dots,h\mod p_r^{a_r})\\
                    &=\varphi(g)\varphi(h)
    \end{align*}

    Injective:

    Let $g \in \ker(\varphi)$. Then
    \[g = 1 \mod p_i^{a_i}\] 
    Since $p_i^{a_i}$ are pairwise coprime, the Chinese remainder theorem gives $g \equiv 1 \mod n$.

    Thus, $g=1$.

    Surjective:

    Let $(g_1, g_2,\dots,g_r) \in \prod (\Z \diagup p_i^{a_i} \Z)^\times$

    By the Chinese remainder theorem, there exists an integer $m$ s.t.
    \[m \equiv g_i \mod \p_i^{a_i}\]
    Thus, $\varphi(m)=(g_1,g_2,\dots,g_r)$.
    Therefore, $\varphi$ is an isomorphism, and
    \[(\Z \diagup n\Z)^\times \cong \prod (\Z \diagup p_i^{a_i} \Z)^\times\]
\end{problemproof}
\end{document}